\Needspace{18\baselineskip}
\subsection{Search coordinates respect renaming}
\label{sec:search-renaming}
Let $a:G\cong H$ be a literal-label isomorphism and write $a_*$ for its port
map. Internal wire labels are preserved by the isomorphism's label field;
boundary labels are preserved because boundary positions are fixed.
\lean{Graph/SearchRenaming}{11}{16}{Linear.iso_vertexLabel}

\begin{definition}[Transport of operation views and coordinates]
Transport an operation view by applying the edge bijection to its edge and
$a_*$ to its source and both ordered lists. Keep its operation symbol.
Source and target certificates follow from the isomorphism's incidence laws;
arity follows because mapping a list preserves length. Transport a typed
coordinate $(m,v)$ to $(m,a_*v)$. Wire-label preservation supplies its typing
certificate.
\end{definition}
\lean{Graph/SearchRenaming}{21}{30}{Linear.SearchOperation.rename}
\lean{Graph/SearchRenaming}{86}{87}{Linear.TypedPosition.rename}

\begin{lemma}[Transport of the local rules]
Operation entry and argument return commute with transport. Every coordinate
search step transports forwards, and a step between transported coordinates
exists if and only if the original step exists.
\end{lemma}
\begin{proof}
Entry depends on whether the argument list is empty or has a first element;
mapping preserves these alternatives. Return depends on the optional element
at index $i+1$; lookup in a mapped list is the mapped lookup. These facts prove
the two local calculation identities.

Induct on the search rule. The variable rule uses the transported label and
singleton source certificate. The operation-entry case uses the entry identity.
For argument return, map the certified lookup at index $i$ and use the return
identity. This proves forward transport for all five rules. Apply the same
result to $a^{-1}$ for reflection: the composite port map is the identity,
and neither map changes the mode. The same inverse law proves that transporting
a typed coordinate by $a$ and back by $a^{-1}$ returns that coordinate; subtype
extensionality discards the irrelevant typing certificates.
\end{proof}
\lean{Graph/SearchRenaming}{32}{35}{Linear.SearchOperation.rename_enter}
\lean{Graph/SearchRenaming}{37}{40}{Linear.SearchOperation.rename_resume}
\lean{Graph/SearchRenaming}{47}{62}{Linear.SearchStep.rename}
\lean{Graph/SearchRenaming}{65}{78}{Linear.SearchStep.rename_iff}
\lean{Graph/SearchRenaming}{90}{95}{Linear.TypedPosition.rename_symm}

\Needspace{18\baselineskip}
\subsection{Insertion respects renaming}
Extend the vertex and edge bijections of $a$ by the identity on each fresh
singleton. The extended port map sends hosted vertices to hosted images and
fixes the fresh vertex. It also commutes with splitting: injectivity gives
$a_*w=a_*v$ exactly when $w=v$, so both splits select the fresh branch in
exactly the same case.
\lean{Graph/InsertionRenaming}{10}{12}{Linear.FocusInsertion.rename_host}
\lean{Graph/InsertionRenaming}{14}{15}{Linear.FocusInsertion.rename_fresh}
\lean{Graph/InsertionRenaming}{17}{25}{Linear.FocusInsertion.rename_split}

\begin{theorem}[Naturality of insertion]
Insertion at $(m,v)$ in $G$ is focus-preservingly isomorphic to insertion at
$(m,a_*v)$ in $H$. In particular, realizing a typed coordinate commutes with
its transport through $a$.
\end{theorem}
\begin{proof}
Use the two extended bijections. Retained labels are preserved by $a$;
fresh vertex and edge labels agree by construction. On retained edges,
commutation with splitting proves the source equation and commutation with
hosting proves the target equation, after applying the corresponding
incidence equation of $a$. The fresh edge has exactly the hosted insertion
position as source and the fresh vertex as target, so the two port identities
prove its equations. The mode is unchanged, and the fresh focus edge maps
to itself. Packaging these fields gives a focussed isomorphism. Applying
this construction to typed coordinates gives the realization result;
focus-free and hole-free certificates ensure both endpoints are states.
\end{proof}
\lean{Graph/InsertionRenaming}{28}{67}{Linear.FocusInsertion.rename_iso}
\lean{Graph/InsertionRenaming}{70}{77}{Linear.FocusInsertion.rename_focussed}
\lean{Graph/InsertionRenaming}{86}{89}{Linear.TypedPosition.realize_rename}

\Needspace{20\baselineskip}
\subsection{Exact characterization of search from realized states}
\label{sec:realized-search}
\begin{lemma}[Recovery of a realized coordinate]
If $p$ is a typed coordinate of a focus-free, hole-free graph $G$, there is
an isomorphism $a$ from the underlying graph of its realization to $G$
such that $a_*\operatorname{coordinate}(\operatorname{realize}(p))=p$.
\end{lemma}
\begin{proof}
Use the position-sensitive insertion--erasure isomorphism. It restores the
vertex, while the extracted mode is definitionally the original mode. Equality
of these two components gives equality of coordinates and then of the typed
subtypes.
\end{proof}
\lean{Graph/RealizedSearch}{14}{22}{Linear.TypedPosition.coordinate_recovery}

\begin{theorem}[Coordinate characterization]
\termdef{coordinate-characterization}
For every actual target state $t$,
\[
 \operatorname{realize}(p)\longrightarrow t
 \quad\Longleftrightarrow\quad
 \exists q,\quad p\longrightarrow_G q\ \land\
       \operatorname{realize}(q)\cong_f t.
\]
Here the left relation is erasure-based state search, the middle relation is
coordinate search on the original graph $G$, and $\cong_f$ preserves the
focus and mode. Both realizations use the same original graph $G$.
\end{theorem}
\begin{proof}
Choose the recovery isomorphism $a$ from the previous lemma. For the forward
direction, a state-search witness consists of a successor coordinate on the
erased graph and an isomorphism from its realization to $t$. Transport this
coordinate and its step through $a$. Coordinate recovery identifies the source
with $p$. Naturality of realization, inverted and composed with the supplied
endpoint isomorphism, identifies the new realization with $t$.

Conversely, transport $q$ and its step through $a^{-1}$. Inverse transport and
coordinate recovery identify the source with the extracted coordinate of
$\operatorname{realize}(p)$. Invert the realization isomorphism and compose it
with the given isomorphism to $t$. The transported coordinate, transported
step, and composite isomorphism are exactly the witnesses required by the
state-search definition.
\end{proof}
\lean{Graph/RealizedSearch}{26}{48}{Linear.TypedPosition.search_realize_iff}
Taking $t=\operatorname{realize}(q)$ and the identity endpoint isomorphism
shows that every coordinate step realizes an actual state-search step between
its two prescribed realizations. This strengthens the earlier endpoint-only
realization theorem. It does not identify search with contextual graph rewriting.
\lean{Graph/RealizedSearch}{52}{55}{Linear.TypedPosition.search_realize}

\Needspace{20\baselineskip}
\subsection{Erasure and state search respect renaming}
\label{sec:erasure-renaming}
Fix a focussed isomorphism $a:f\cong_f g$. For any bijection carrying $x$ to
$y$, restrict it to the complements of $x$ and $y$. Its forward and inverse
maps are the original maps with non-deletion certificates: injectivity prevents
a retained point from mapping to the deleted point, and the inverse law proves
the analogous claim backwards. Subtype extensionality lifts both inverse laws.
\lean{Graph/ErasureRenaming}{10}{15}{Linear.exceptOfEq}

The graph vertex bijection carries the selected internal target of $f$ to
that of $g$. Indeed, the focussed isomorphism identifies the targets, their
selection certificates make them internal, and the internal constructor is
injective. Restrict this vertex bijection at the selected targets and the edge
bijection at the focus edges. These give bijections of the erased carriers.
\lean{Graph/ErasureRenaming}{21}{24}{Linear.Focussed.Iso.targetVertex}
\lean{Graph/ErasureRenaming}{27}{28}{Linear.Focussed.Iso.erasedVertices}
\lean{Graph/ErasureRenaming}{31}{32}{Linear.Focussed.Iso.erasedEdges}

\begin{lemma}[Naturality of contraction]
Inclusion of retained ports commutes with the restricted vertex bijection.
The joined positions correspond, and contraction commutes with renaming.
\end{lemma}
\begin{proof}
Commutation with inclusion is a case analysis on ports. To compare joined
positions, apply the injective inclusion on the target side. The resulting
equation is precisely preservation of the focus source. For contraction,
separate the deleted target from other vertices. In the first case both sides
are the corresponding joined positions. In the second, injectivity ensures
that the image is not the target either. Apply inclusion on the target side:
both contraction--inclusion composites reduce to their original vertices,
and the inclusion commutation identity proves equality.
\end{proof}
\lean{Graph/ErasureRenaming}{34}{37}{Linear.Focussed.Iso.erased_lift}
\lean{Graph/ErasureRenaming}{40}{45}{Linear.Focussed.Iso.erased_position}
\lean{Graph/ErasureRenaming}{48}{63}{Linear.Focussed.Iso.erased_contract}

\begin{theorem}[Naturality of erasure and state search]
The restricted bijections form an isomorphism of the erased graphs, carrying
the extracted coordinate of the first state to that of the second. Search
is invariant under focus-preserving isomorphisms in both arguments.
\end{theorem}
\begin{proof}
Retained labels inherit the original isomorphism's label equations. For each
incidence list, use its original equation and commute contraction past the
port map pointwise. This proves the erased graph isomorphism. Its coordinate
equation follows from preservation of the mode and the joined position.

Given a search step from the first state, transport its successor coordinate
through the erased isomorphism. The coordinate equation identifies the new
source, and search transport supplies the step. Invert the realization
isomorphism and compose with the old endpoint isomorphism to retain the same
actual target. This proves source invariance. Combine it with the earlier
target invariance to obtain invariance in both arguments; use the inverse
isomorphisms for the reverse implication.
\end{proof}
\lean{Graph/ErasureRenaming}{66}{86}{Linear.Focussed.Iso.erased}
\lean{Graph/ErasureRenaming}{94}{99}{Linear.State.coordinate_iso}
\lean{Graph/ErasureRenaming}{102}{112}{Linear.State.search_source_iso}
\lean{Graph/ErasureRenaming}{115}{122}{Linear.State.search_iso_iff}

\Needspace{20\baselineskip}
\subsection{Search on isomorphism classes of states}
\label{sec:state-quotient}
\begin{definition}[State isomorphism class]
\termdef{state-isomorphism-class}
Two states are equivalent when there exists a focus-preserving literal-label
isomorphism between them. Reflexivity, symmetry, and transitivity follow
from the corresponding isomorphism constructions. A state class is the
quotient by this equivalence; write $[s]$ for the class of $s$.
\end{definition}
This quotient preserves full edge labels literally, including box bodies.
It does not identify independently renamed contents of nested boxes.
Section~\ref{sec:recursive-state-search} constructs the recursive state quotient
and proves projection and exact-length search lifting between the two quotients.
\lean{Graph/StateQuotient}{10}{13}{Linear.stateSetoid}
\lean{Graph/StateQuotient}{16}{16}{Linear.StateClass}
\lean{Graph/StateQuotient}{22}{22}{Linear.StateClass.ofState}

\begin{definition}[Quotient search and observations]
\termdef{quotient-search}
Define $[s]\longrightarrow[t]$ by state search $s\longrightarrow t$.
The preceding two-argument invariance proves independence of representatives,
so quotient elimination defines a proposition on classes. Initiality and
finality descend in the same way, using their focussed-isomorphism invariance.
\end{definition}
On supplied representatives, all three descended propositions agree with the
original propositions by definition. The representative equations are therefore
reflexive equivalences, not extra assumptions.
\lean{Graph/StateQuotient}{25}{27}{Linear.StateClass.Search}
\lean{Graph/StateQuotient}{29}{30}{Linear.StateClass.search_iff}
\lean{Graph/StateQuotient}{39}{41}{Linear.StateClass.Initial}
\lean{Graph/StateQuotient}{44}{46}{Linear.StateClass.Final}
\lean{Graph/StateQuotient}{48}{48}{Linear.StateClass.initial_iff}
\lean{Graph/StateQuotient}{50}{50}{Linear.StateClass.final_iff}

\begin{theorem}[Deterministic quotient search and its boundaries]
If $S\longrightarrow T$ and $S\longrightarrow U$, then $T=U$.
A final class has no search successor, and no search step reaches an initial
class.
\end{theorem}
\begin{proof}
Eliminate the three quotient arguments to representatives $s,t,u$.
The two representative steps have a common source, so state-level determinism
gives a focussed isomorphism $t\cong_f u$. Quotient soundness identifies their
classes. For each boundary theorem eliminate its two class arguments;
the goal becomes the already proved state-level boundary law.
\end{proof}
\lean{Graph/StateQuotient}{33}{36}{Linear.StateClass.search_deterministic}
\lean{Graph/StateQuotient}{53}{56}{Linear.StateClass.final_no_search}
\lean{Graph/StateQuotient}{59}{62}{Linear.StateClass.no_search_initial}
These are results for the search relation alone. Substitution, computation,
rootedness and refocusing remain necessary for
the full concrete UAM argument.

\Needspace{18\baselineskip}
\subsection{Concrete quotient-search regression}
The passive nullary value step now runs from the previously defined
\texttt{initialValue} to the specific \texttt{finalValue}, using the general
coordinate-to-state theorem. Erase and reconstruct \texttt{initialValue} to
define \texttt{recoveredValue}. Its carriers are rebuilt by deletion and
insertion, but state recovery supplies a focussed isomorphism to the original.
Source invariance therefore gives a step from this reconstructed state to
\texttt{finalValue}. Quotient soundness identifies its class with the original
initial class. Finally, quotient determinism compares any successor of that
initial class with the concrete value step: it must equal the class of
\texttt{finalValue}.
\lean{Graph/QuotientSearchExamples}{9}{12}{Linear.Examples.value_search_final}
\lean{Graph/QuotientSearchExamples}{15}{17}{Linear.Examples.recoveredValue}
\lean{Graph/QuotientSearchExamples}{19}{20}{Linear.Examples.recovered_value_search}
\lean{Graph/QuotientSearchExamples}{22}{23}{Linear.Examples.recovered_value_class}
\lean{Graph/QuotientSearchExamples}{26}{29}{Linear.Examples.value_class_successor}
